\documentclass[12pt]{article}

\usepackage{amsmath, amssymb, amsthm}
\usepackage{algorithm}
\usepackage{algorithmic}
\usepackage{geometry}
\usepackage{enumitem}
\usepackage[hidelinks]{hyperref}   % <-- This removes red boxes

\geometry{margin=1in}

\title{\textbf{Lecture 10: Randomized Min-Cut Algorithm}}
\author{}
\date{}

\begin{document}

\maketitle
\tableofcontents
\newpage

%------------------------------------------------
\section{Overview}

\subsection{Min-Cut Problem}
The Min-Cut problem asks to find the smallest set of edges whose removal disconnects a graph.

\subsection{Why Use Min-Cut?}
Min-cut helps identify weakest connectivity in a network.

\subsection{What is Min-Cut?}
Given a graph $G=(V,E)$ with $n$ vertices, a minimum cut is a cut-set of minimum cardinality.

\subsection{Successful Min-Cut Run}
The algorithm correctly identifies the minimum cut of the graph.

\subsection{Unsuccessful Min-Cut Run}
The algorithm contracts critical edges early and fails to output the true minimum cut.

\subsection{Max-Flow Min-Cut Theorem}
In a flow network, the maximum flow from source $S$ to sink $T$ equals the capacity of the minimum cut.

\subsection{Deterministic Min-Cut Algorithm}
Guarantees exact solution with higher time complexity.

\subsection{Stoer--Wagner Minimum Cut Algorithm}
Deterministic polynomial-time algorithm for global minimum cut.

\subsection{Randomized Min-Cut Algorithm}
Uses random edge contractions to find minimum cut probabilistically.

\subsection{Why Randomized Algorithm?}
\begin{itemize}
    \item Probabilistic guarantees
    \item Faster for large graphs
    \item Avoids worst-case inputs
\end{itemize}

%------------------------------------------------
\section{Definitions and Notation}

\subsection{Cut-Set}
A cut-set is a set of edges whose removal disconnects the graph.

\subsection{Min-Cut Problem}
Given $G=(V,E)$ with $n$ vertices, find a minimum cardinality cut-set.

\subsection{Edge Contraction}
For edge $(u,v)$:
\begin{itemize}
    \item Merge $u$ and $v$
    \item Remove self-loops
    \item Keep parallel edges
\end{itemize}

\subsection{Flow Network Terms}
\begin{itemize}
    \item Capacity of a cut = Sum of capacities from $X$ to $Y$
    \item Minimum cut = Cut with smallest capacity
    \item Maximum flow = Largest flow from $S$ to $T$
\end{itemize}

%------------------------------------------------
\section{Assumptions and Conditions}

\begin{itemize}
    \item Graph $G=(V,E)$ is given
    \item Treated as undirected multigraph
    \item Randomized algorithms provide probabilistic guarantees
\end{itemize}

%------------------------------------------------
\section{Main Results / Theorems}

\subsection{Max-Flow Min-Cut Theorem}
The maximum flow equals the minimum cut capacity.

\subsection{Stoer--Wagner Theorem}
Let $s$ and $t$ be vertices of $G$. Let $G/\{s,t\}$ be graph after merging $s$ and $t$.

Minimum cut of $G$ is the smaller of:
\begin{itemize}
    \item Minimum $s$--$t$ cut of $G$
    \item Minimum cut of $G/\{s,t\}$
\end{itemize}

\subsection{Theorem for Min-Cut Set}
Karger's algorithm outputs a minimum cut with probability at least:
\[
\frac{2}{n(n-1)}
\]

%------------------------------------------------
\section{Proof / Justification}

\subsection{Stoer--Wagner Justification}
Either:
\begin{itemize}
    \item Minimum cut separates $s$ and $t$
    \item Or contraction preserves minimum cut
\end{itemize}

%------------------------------------------------
\section{Deterministic vs Randomized}

\subsection{Deterministic (Stoer--Wagner)}
Time Complexity:
\[
O(VE + V^2 \log V)
\]

\subsection{Randomized (Karger)}
Time Complexity:
\[
O(V^2)
\]

%------------------------------------------------
\section{Pseudocode}

\begin{algorithm}
\caption{MinimumCutPhase(G, a)}
\begin{algorithmic}
\STATE $A \leftarrow \{a\}$
\WHILE{$A \neq V$}
    \STATE Add most tightly connected vertex to $A$
\ENDWHILE
\STATE Return cut-of-the-phase
\end{algorithmic}
\end{algorithm}

\begin{algorithm}
\caption{MinimumCut(G)}
\begin{algorithmic}
\WHILE{$|V| \ge 1$}
    \STATE Choose any $a \in V$
    \STATE MinimumCutPhase(G, a)
    \STATE Update minimum cut if lighter
    \STATE Merge last two added vertices
\ENDWHILE
\STATE Return minimum cut
\end{algorithmic}
\end{algorithm}

%------------------------------------------------
\section{Python Simulation}

\begin{verbatim}
import random

def karger(graph):
    while len(graph) > 2:
        u = random.choice(list(graph.keys()))
        v = random.choice(graph[u])
        graph[u].extend(graph[v])
        for vertex in graph[v]:
            graph[vertex].remove(v)
            graph[vertex].append(u)
        graph[u] = [x for x in graph[u] if x != u]
        del graph[v]
    return len(list(graph.values())[0])
\end{verbatim}

\end{document}
